% !TEX root = ../thesis.tex
\chapter{Developing an Algorithm}

Now that we've formally defined our problem and ideal solution, we need to
actually create a way to find that solution. We also need to keep in
mind that any algorithm we develop must improve upon the current system, as
there already exists a standard method for HUM 110 sign-ups. So, we will evaluate the current
system against our newfound ISF1M goals and determine if we can develop an
algorithm that accomplishes those goals more effectively.

\section{First Come, First Served}
As it currently stands, the sign-up process for HUM 110 classes takes place during
freshman class registration. The registration software, SOLAR, opens to incoming
freshmen at a certain time, and at that time (or later) all freshmen log on to
the platform and choose their own academic year schedule in whatever order they
desire. This includes their HUM 110 section (or, as is the driving force behind this
project, sections plural). Though this process may not inherently seem like an
algorithm, it can actually be modelled as the algorithm commonly known as
\textit{First Come First Served} (\textsc{fcfs}).

\textsc{fcfs} takes some ordering of students that represents the registration
order (in this case, those who sign up for HUM 110 the fastest to slowest after
the registration software opens), and at their point of
selection each student chooses their top available choice until every student
has chosen a HUM class. In this case, “available” means classes that have not
yet reached capacity. We write out \textsc{fcfs} formally in Algorithm
\ref{alg:fcfs}, where
$P_r$ represents the ordered preference list for a student $r$ (corresponding to
their preference relation defined in Definition \ref{def:overview}).

\begin{algorithm}
  \caption{First Come First Served}
  \label{alg:fcfs}
  \begin{algorithmic}[1]
    \For{each student $r$ in $\mathcal{R}$} \Comment{where $\mathcal{R}$ is an ordered set}
    \For{each class $C$ in $P_r$}  
    \If{\# students assigned to $C < q$}
    \Comment{class is not at capacity} 
    \State $\mu(r) \gets C$ 
    \State \textbf{break}  \Comment{move to next student} \EndIf \EndFor \EndFor
    \State \Return matching $\mu$

  \end{algorithmic}
\end{algorithm}

There are various benefits to a \textsc{fcfs} algorithm, one of the most
significant for our purposes being that the given matching will be
\textit{inter-group stable}. Because each student chooses their most preferred
option when it is their turn to select, there can be no two pairs of students
that wish to switch. Theorem \ref{thm:fcfs-stable} proves \textsc{fcfs}
stability later in this chapter. Additionally, in practice many students get
their first or second choice of HUM section under \textsc{fcfs}, and the
algorithm is easy to implement because it is essentially already in use.

However, \textsc{fcfs} has other issues. For instance, it is possible some
students will not get to choose a class they want at all if their preferred
sections are already at capacity. Self-selection also allows for inconsistent class sizes
and it does not consider student groups, thus permitting an inconsistent
distribution of students. Our goal in this section is to develop an algorithm
that improves upon \textsc{fcfs} with group fairness.

\section{A Greedy Approach}

We introduce group fairness into our algorithm with what we call a
\textit{greedy} algorithm: one that attempts to maximize the rank a student
gives their assignment while following the group fairness constraints. We call it \textit{greedy}
because for each student, we still assign them their best available choice at
the time of evaluation. However, in addition to the class capacity constraint,
“available” also considers the group caps on the seats from Definition
\ref{def:group-caps}. So, a class that is not yet at
capacity is still unavailable to a student of group $T$ if it has already
reached its maximum number of group $T$ students. Essentially, \textsc{greedy}
implements \textsc{fcfs} under the additional group fairness restriction.

\begin{algorithm}
  \caption{Greedy}
  \label{alg:greedy}
  \begin{algorithmic}[1] \For{each student $t$ in $\mathcal{R}$} \Comment{where $t$ is a
    student of some group $T \subset \mathcal{R}$} \For{each class $C$ in $P_t$}
    \If{\# students assigned to $C
        < q$ and \# $T$ students assigned to $C < q^T$} \State $\mu(t) \gets C$
        \State \textbf{break}  \Comment{move to next student} \EndIf \EndFor
        \EndFor
    \State \Return matching $\mu$
  \end{algorithmic}
\end{algorithm}

Let's compare \textsc{greedy} to \textsc{fcfs}. The primary difference is in
line 3, where we add the group cap constraint as part of the conditional check
for whether a class is available. Recall from Chapter 1 that a ``vanilla"
environment is the most basic version of a problem; in our case, a vanilla
environment removes group fairness constraints, simplifing the problem to
treat all students the same when assigning class sections. We can start by showing that in a
vanilla environment, both algorithms output a stable matching. 

\begin{lemma}
  \label{thm:fcfs-greedy}
  \textsc{fcfs} is equivalent to \textsc{greedy} in a vanilla environment.
\end{lemma}

\begin{proof}
  Observe that in a vanilla environment, $q^T = q$ because there are
  no group caps. So, the second part of the \texttt{and} statement in line 3 of
  Algorithm \ref{alg:greedy}
  will always be true when the first part of the statement is true. Therefore,
  both algorithms will behave equivalently on the same ordered input of students
  in the first \texttt{for} loop, as the conditional statements check the same condition
  as line 3 of Algorithm \ref{alg:fcfs}. Thus, \textsc{fcfs} is equivalent to
  \textsc{greedy} in a vanilla environment.
\end{proof}

\begin{lemma}
  \label{thm:greedyvanilla}
  \textsc{greedy} outputs an inter-group stable matching in a vanilla environment.
\end{lemma}

\begin{proof}
  Let $\mu$ denote the outputted matching over an ordered set of
  students, $\mathcal{R}$. Let $r_1, r_2 \in \mathcal{R}$, where $r_1$ comes
  before $r_2$ in the order.
  Assume for the sake of a contradiction that $r_1$ and $r_2$ form a blocking
  pair, meaning $\mu(r_2) >_{r_1} \mu(r_1)$ and $\mu(r_1) >_{r_2} \mu(r_2)$.

  By line 3 of Algorithm \ref{alg:greedy}, we know $\mu(r_1)$ was either $r_1$'s top
  available choice or $r_1$ is unmatched. So, if $\mu(r_2) >_{r_1} \mu(r_1)$, then
  $\mu(r_2)$ must not be available when $r_1$ is matched, which means $r_2$ was
  matched first. However, this is a contradiction of our assumption that $r_1$
  comes before $r_2$ in the input to the algorithm. Thus, $\mu(r_2)$ is available
  when $r_1$ is matched, so it must be the case that $\mu(r_1) >_{r_1} \mu(r_2)$.
  Therefore, $r_1$ and $r_2$ do not form a blocking pair, and we can conclude that
  $\mu$ is stable.
\end{proof}

\begin{theorem}
  \label{thm:fcfs-stable}
  \textsc{fcfs} outputs an inter-group stable matching.
\end{theorem}

\begin{proof}
  By Lemma \ref{thm:fcfs-greedy}, we know \textsc{fcfs} is equivalent to
  \textsc{greedy} in a vanilla environment. Lemma \ref{thm:greedyvanilla} proves
  that in a vanilla environment, \textsc{greedy} outputs an inter-group stable
  matching. So, together those two lemmas imply that \textsc{fcfs} outputs an
  inter-group stable matching.
\end{proof}

Now, we can formally evaluate what happens when we do consider groups. As we
discussed in Theorem \ref{thm:unstable}, we know an inter-group stable matching
may not exist at all when we introduce caps on groups. However, even assuming there does
exist an inter-group stable matching, \textsc{greedy} is still not guaranteed to find it.

\newpage 

\begin{theorem}
  \label{thm:greedyunstable}
  \textsc{greedy} does not always output a inter-group stable matching in a
  general environment.
\end{theorem}

\begin{proof}
  Let $\mathcal{R} = X \cup Y$ where there are two groups of students, group $X$
  and group $Y$. Let  $\mathcal{C} = \{ C_1, C_2, C_3 \mid q^X = 1, q^Y =
  1$\}, meaning each $C \in \mathcal{C}$ can have at most one student from each
  group. Let $\mathcal{R} = \{x_1, y_1, x_2, y_2\}$ be ordered where students $x_1, x_2 \in X$ and
  $y_1, y_2 \in Y$. Let the preference lists of the students be as follows: \[
    \begin{array}{ll} x_1: & C_1 \succ C_3 \\ y_1: &
             C_2 \succ C_3             \\ x_2: & C_1 \succ C_2\\ y_2: & C_2 \succ
                C_1.
    \end{array}
  \] 

  Observe that an inter-group stable matching $\mu'$ does exist in this
  environment: \[ 
    \begin{array}{ll}
      x_1 &\gets C_1 \\ 
      y_1 &\gets C_3 \\
      x_2 &\gets C_2 \\
      y_2 &\gets C_2. 
    \end{array}
\]

  We know $\mu'$ is stable because $x_1$ and $y_2$ both received their top
  choice, so neither of them wish to swap with another student. $y_1$
  wishes to swap with $x_2$, but $\mu'(x_2) \succ_{x_2} \mu'(y_1)$, so they do
  not form a blocking pair. Similar logic follows for $y_2$, who $y_1$ would
  also want to swap with. $x_2$ wants to swap with $x_1$, who we already know
  does not wish to swap with anyone. Therefore, $\mu'$ does not have any
  blocking pairs, and thus it is inter-group stable.

  Now let $\mu$ denote \textsc{greedy}'s outputted matching over a set of
  students, $\mathcal{R}$. We will prove by counter example that $\mu$ is not
  stable. Observe that $\mu(x_1) = s_{c_1}^X$ and $\mu(y_1) = s_{c_2}^Y$:
  because they are the first two students to be matched and their first choices
  are different, they both receive their top choices. Next, $\mu(x_2) = s_{c_2}^X$
  because $C_1$ already has a student in group $X$, so $C_1$ is no longer
  available when matching $x_2$, and similarly $\mu(y_2) = s_{c_1}^Y$ because
  $C_2$ already has a student in group $Y$. Notice that $\mu(y_2) >_{x_2}
  \mu(x_2)$ and $\mu(x_2) >_{y_2} \mu(y_2)$, which forms a blocking pair where
  $x_2$ and $y_2$ are incentivized to switch classes with each other. So, it
  must be the case that $\mu$ is not an inter-group stable matching and
  therefore \textsc{greedy} cannot guarantee an inter-group stable matching
  under group cap constraints, even when an inter-group stable matching exists.
\end{proof}

Notably, the above counterexample does exhibit intra-group stability as
$\mu(x_1)$ is $x_1$'s top choice, so they will never want to swap with $x_2$
(and similar for $y_1$ and $y_2$). In fact, we claim that \textsc{greedy} will
always output an intra-group stable matching.
At each point of evaluation for a student of some group, they are given their
highest ranked available seat for that group. So within groups, we have a
similar stability argument to Lemma \ref{thm:greedyvanilla}.

While \textsc{greedy} is intra-group stable, its primary issue is that
it does not guarantee a full matching. Students can be left unmatched if every
class section on their preference list is already at capacity for their gorup.
We need to come up with an algorithm that makes use of our
helpful graph framework from Definition \ref{weighted-def} if we want to get
closer to creating ISF1M solutions.

\section{The Weight of the Situation}
Let's return to our conception of the HUM-assignment problem as modeled by a graph. We
know by Theorem \ref{max-weight} we can use a maximum-weight matching algorithm
to find a solution that is full and group fair. Notably, this is only one
requirement away from being an ISF1M solution
(see Definition \ref{def:ispim}): we just need to reintroduce intra-group
stability to our max-weight matching. To do that, we need to
remove all blocking pairs between students of the same group while preserving
the weight of the max-weight matching, as the weight is how we determined that
the matching is full and group fair. Here we define Algorithm \ref{alg:wt}, called
the \textit{Weighted-Trading} (\textsc{wt}) algorithm.

\begin{algorithm}
  \caption{Weighted-Trading}
  \label{alg:wt}
  \begin{algorithmic}[1] \State $w(\mu) \gets 0$ \While{$w(\mu)
    \neq |\mathcal{R}| + |\mathcal{S}_{res}|$ and \# reserved seats for all groups $> 0$}
    \State $\mu \gets \textsc{max-weight}(\mathcal{R}, \mathcal{S})$ \If{$w(\mu)
    \neq |\mathcal{R}| + |\mathcal{S}_{res}|$} \For {each group of seats $T
  \subset \mathcal{S}$} \State $|T| =
  |T| - 1$ \Comment{reduce the size of $\mathcal{S}_{res}$} \EndFor
  \EndIf \EndWhile \For{each group $T \subset \mathcal{R}$}
    \While{$\exists$ blocking pair $t_i, t_j$} \State $\mu(t_i) \gets \mu(t_j)$,
    $\mu(t_j) \gets \mu(t_i)$ \EndWhile \EndFor \State \Return matching $\mu$
  \end{algorithmic}
\end{algorithm}

We can break the \textsc{wt} algorithm into two key stages: lines 1 to 9 and
lines 10 to 15.

The first stage is the \textit{weighted} part of the algorithm: it repeatedly
runs \textsc{max-weight} until it finds a satisfying assignment. If the weight of
the first \textsc{max-weight} run is not $|\mathcal{R}| + |\mathcal{S}_{res}|$ (as per Theorem
\ref{max-weight}), we ``relax" our seat reservations to allow for more students
of various groups to fill the class. This involves decrementing the number of
reserved seats per group by one and then re-running \textsc{max-weight} now that
the set of reserved seats is smaller. If the seat reservations are relaxed to
the point that there are no more reserved seats, it means there is no group-fair
solution and we simply use the \textsc{max-weight} output to get a full matching for
students without considering group caps. Then, we still proceed with the trading
stage for intra-group stability.

The latter stage is the \textit{trading} part of the algorithm, which handles
our intra-group stability requirement. We check for all existing blocking pairs
between students of the same group, and we swap or trade their assignments with
each other. Notably, we have to continually check each pair in a group of
students until no more blocking pairs remain, as checking each student once does
not account for new blocking pairs that may arise from swapping the assignments
of two students. Example \ref{eg:trading} demonstrates this complication. At the
end of the algorithm, we have a full, group fair, intra-group
stable matching (or an ISF1M) where each student is assigned a class they prefer
if such a matching exists.

\begin{example} \label{eg:trading}
  Let $x_1, x_2, x_3 \in X \subset \mathcal{R}$ and $\mathcal{C} = \{C_1, C_2, C_3\}$.
  Let the preference lists of the students be as follows: \[
    \begin{array}{ll} x_1: & C_1 \succ C_2 \succ C_3 \\ x_2: &
             C_3 \succ C_2             \\ x_3: & C_2 \succ C_1.    \end{array}
  \]
  Let the $\mu$ be the max-weight matching output at line 9 of Algorithm
  \ref{alg:wt} where  \[
    \begin{array}{ll} \mu(x_1) &= C_3 \\ \mu(x_2) &= C_2\\ \mu(x_3) &= C_1.    \end{array}
  \]

  In this case,  $\mu(x_2) >_{x_1} \mu(x_1)$ and $\mu(x_1) >_{x_2} \mu(x_2)$,
  where $x_1$ and $x_2$ form the only blocking pair. In the trading stage, we swap their
  assignments so that $\mu(x_1) = C_2$ and $\mu(x_2) = C_3$. Now, observe that
  $x_1$ and $x_3$ form a new blocking pair, though they didn't before. When we make
  the swap between $x_1$ and $x_2$, it causes $\mu(x_1) >_{x_3} \mu(x_3)$,
  while $\mu(x_3) >_{x_1} \mu(x_1)$ was already true. So, when performing
  trading, we have to continually check blocking pairs for each student because
  new ones may arise during the process. 
\end{example}

We can assert that the outputted matching after the trading stage is still full and group
fair by proving that the weight of the graph is still the same maximum weight found
in the weighted stage of \textsc{wt}.

\begin{theorem}
  \label{thm:same_weights}
  For a matching $\mu$ output by the \textsc{wt} algorithm, $w(\mu) =
  |\mathcal{R}| + |\mathcal{S}_{res}|$ if such a matching exists.
\end{theorem}

\begin{proof}
  At line 9 of Algorithm \ref{alg:wt}, $w(\mu) = |\mathcal{R}| + |\mathcal{S}_{res}|$
  because it was just output by \textsc{max-weight}, assuming such a
  matching exists. So, we have to
  evaluate the \texttt{for} loop that compares each student of one group to each other
  student of that same group. In the \texttt{for} loop, the students swap seats if they
  both prefer each other's match in $\mu$. This means that each student still
  prefers their new match, so the weight of the edge retains weight at least 1
  from that student. Then, if either student is in a seat reserved for that
  group, we replace them with a student of the same group, so the edge retains the
  additional weight of the reserved seat. Therefore, swapping any two students of
  the same group does not change the weight of $\mu$.
\end{proof}

While it would be nice to guarantee inter-group stability, recall that the only
known algorithms for assignment problems that consider inter-group stability and
group fairness are \textit{NP-complete}  \cite{santhini}. With the \textsc{wt} algorithm, we can guarantee that if an ISF1M exists, we can find it in an efficient amount of time.

\begin{theorem}
  \label{thm:polywt}
  The \textsc{wt} algorithm will output an intra-group stable matching in cubic time with
  respect to the vertices in $G$.
\end{theorem}

\begin{proof}
  Let $\mu$ be the matching output by \textsc{max-weight} at
  line 9 of Algorithm \ref{alg:wt}. Observe lines 11 and 12, which outline that
  within each student group, we must swap the
  assignments of blocking pairs until no more remain. We know this loop will
  terminate because of our assumption that preference lists are finite, so even
  in the worst case
  students will eventually reach their highest possible preference and no longer
  be able to form blocking pairs. The maximum length of a preference list is
  equivalent to the number of classes, $|\mathcal{C}|$. So, the maximum number of times
  the \texttt{while} loop can execute per \texttt{for} loop is as follows: \[
  |\mathcal{C}| \times |\mathcal{R}|.\]
  This is because at most we can traverse each student's entire preference list once. If
  we terminate this loop for each group of student $T \subset \mathcal{R}$, then we must have
  intra-group stability because no more blocking pairs remain in each group.

  We can analyze the runtime of this algorithm. Since the \texttt{while} loop executes at
  most $|\mathcal{C}| \times |\mathcal{R}| \times \# \text{ groups}$ times, we can bound the
  execution of the \texttt{while} loop to $O(|\mathcal{C}| \times |\mathcal{R}| \times \# \text{
  groups})$. Notably, $|\mathcal{C}| \leq |\mathcal{R}|$ and $\# \text{ groups }
  \leq |\mathcal{R}|$. In practical cases there will be significantly fewer
  student groups and class sections than there are students, so the significant
  factor for this bound is $|\mathcal{R}|$. Therefore, the trading stage of the algorithm
  will run with respect to the size of $\mathcal{R}$.

  Observe that before that, we run \textsc{max-weight} at least once, which runs in $O(V^3)$
  time where $V$ is the set of vertices \cite{hungarian}. Our vertices are
  all the students and seats, so the runtime for the weighted stage is
  $(|\mathcal{S}| + |\mathcal{R}|)^3$. So, our total runtime for the algorithm
  is $O((|\mathcal{S}|+ |\mathcal{R}|)^3 + (|\mathcal{C}| \times |\mathcal{R}|
  \times \# \text{ groups}))$ = O$((|\mathcal{S}| + |\mathcal{R}|)^3)$. Therefore, the \textsc{wt}
  algorithm outputs an intra-group stable matching in cubic time with respect to
  the vertices in $G$.
\end{proof}

By Theorem \ref{thm:polywt}, we now have an algorithm that theoretically gives
us a group-fair, intra-group stable matching. Note that the matching is also
full, as every student is matched to a class in the \textsc{max-weight} section
if such a matching is possible. Thus, the \textsc{wt} algorithm we've developed
outputs an ISF1M solution to the HUM-assignment problem when one exists.
The next step is to evaluate it against \textsc{fcfs} and \textsc{greedy} to
determine whether it performs well in practice.

